%!TEX root = ./../../springer-mono-format/book.tex
%!TEX root = ./../../thesis-trier.tex

In \cref{sync_max:grp:prim:sec:min_aut,sync_max:perm:sec:aut_sync_words}
we characterized primitivity by extremal state complexity properties of associated automata.
Here, we discuss some more relations
to the fields of state and descriptional complexity;
see~\cite{GaoMRY17,HolzerK11} for surveys of both fields.
In particular, we discuss so-called most complex families of languages. 
Then we discuss the relation of primitivity
to $A$-primitivity
and give an overview on some research
concerned with transformation semigroups
generated by transformations
having certain properties. Lastly, we give some historical remarks on primitive permutation groups.

Synchronizing groups were introduced in~\cite{ArnoldS06}. Since then, they have been studied extensively, including generalizations such as almost synchronizing groups~\cite{DBLP:journals/tcs/AraujoBC13} and connections with graph theory~\cite{doi:10.1112/plms/pdw040,ArCamSt2017}. For a survey, see~\cite{ArCamSt2017}.

In~\cite{Volkov24}, it was shown that automata admitting no non-trivial congruences and containing a unimodal letter are synchronizing. This result is more general than~\cref{syncmax:cor:unimodal-sync}, since such an automaton need not have any permutational letters at all, see also the discussion in~\cite{Volkov24}.

In \cref{sync_max:perm:sec:aut_sync_words} we characterized primitive groups by
the complete reachability of associated automata and transformation semigroups
generated by adjoining a single transformation of defect one.
Conditions that yield the complete reachability 
of transformation semigroups generated by
an imprimitive group and a single transformation of defect one have been investigated by Casas~\cite{CasasTorres2024CompletelyReachable}.

\IN{most complex regular languages}
\paragraph{Most Complex Languages.} \Cref{sync_max:grp:prim:sec:min_aut} and~\ref{sync_max:perm:sec:aut_sync_words} linked the absence
of non-trivial $G$-congruences to
extremal properties of the state complexity of $\Syn(\mathcal A)$, or
of languages accepted by $\mathcal A$, for automata $\mathcal A$
with $G \subseteq \TS{\mathcal A}$.
Restivo \& Vaglica~\cite[Prop.~1]{DBLP:journals/tcs/RestivoV12} (see also Davies~\cite[Prop.~3.2]{DBLP:journals/mst/Davies18} or~\cite[Prop.~4.3.2]{DBLP:phd/basesearch/Davies19})
firstly noted that minimality of a strongly connected automaton
for each distribution of a start state 
and a proper non-empty set of final states
is equivalent to absence of non-trivial $\TS{\mathcal A}$-congruences.
Can we relate other extremal properties investigated in the field of state complexity to $\TS{\mathcal A}$-congruences?
Brzozowski~\cite{DBLP:journals/ijfcs/Brzozowski13,DBLP:conf/dcfs/BrzozowskiD18,DBLP:journals/dmtcs/BrzozowskiDL16}
called a sequence of regular languages $L_n \in \mathcal C$
\emph{most complex in a class $\mathcal C$ of regular languages}  with $n \ge k$ for some $k \ge 1$
if the following conditions hold (for definitions of atoms, reversal, etc., we refer to~\cite{DBLP:journals/ijfcs/Brzozowski13}):
\begin{enumerate}
	\item The minimal automaton for $L_n$ has $n$ states.
	\item The transformation monoid of the minimal automaton has maximal cardinality for all languages in the class $\mathcal C$. 
	\item The minimal automata of the left quotients have $n$ states.
	\item The number of distinct atoms (intersections of complemented and uncomplemented left quotients) of $L_n$ is maximal possible
	and each atom has maximal state complexity for languages in the class $\mathcal C$.
	\item The resulting minimal automata for the following operations are maximal possible w.r.t.~the class $\mathcal C$: reversal,
	star of $L_n$, and product (concatenation), union, intersection, set difference and symmetric difference of $L_m$
	and $L_n$.
\end{enumerate}

For the class of all languages, in~\cite{DBLP:journals/ijfcs/Brzozowski13} the sequence of languages accepted
by the automata $\automat{U_n}{[n]}{\{a,b,c\}}{\delta_n}{\{0\}}{\{n-1\}}$
shown in \cref{syncmax:fig:most-complex-all-lang} was given
where $a$ and $b$ induce $(0 ~ 1 \cdots n - 1)$,
and $(0 ~ 1)$, respectively, on the states
and $c$ is a simple idempotent letter mapping $n-1$ to $0$.
Furthermore, the following subclasses~$\mathcal C$ have
been investigated: union-free regular languages~\cite{DBLP:conf/dcfs/BrzozowskiD18},
ideal languages~\cite{DBLP:journals/dmtcs/BrzozowskiDL16}, and non-returning regular languages~\cite{brzozowski2019most}.
The most complex sequences of languages from~\cite{DBLP:journals/ijfcs/Brzozowski13,DBLP:conf/dcfs/BrzozowskiD18,DBLP:journals/dmtcs/BrzozowskiDL16} 
are all over fixed alphabets, whereas the sequence from~\cite{brzozowski2019most} is over a growing alphabet.

\begin{figure}[ht]
	\centering 	
	\begin{tikzpicture}
		[node distance=5em and 6.5em, on grid, auto, state/.style = {draw = black, rounded rectangle, minimum width=4em, minimum height=1em},every loop/.style={looseness=10},]
		
		\node[state, initial] (0)     {$0$};
		\node[state] (1) [right=of 0] {$1$};
		\node[state] (2) [right=of 1] {$2$};
		\node (dots) [right=of 2] {$\ldots$};
		\node[state] (n2) [right=of dots] {$n-2$};
		\node[state, accepting] (n1) [right=of n2] {$n-1$};
		
		\path[>=stealth] 
		(0) edge node {$a,b$} (1)
		(1) edge node {$a$}   (2)
		(2) edge node {$a$}   (dots)
		(dots) edge node {$a$} (n2)
		(n2) edge node {$a$} (n1)
		(n1) edge [bend left=16, below] node {$a,c$} (0);
		
		\path[>=stealth] 
		(1) edge [bend right=60,above] node {$b$} (0)
		(0) edge [loop above] node {$c$} (0)
		(1) edge [loop above] node {$c$} (1)
		(2) edge [loop above] node {$c$} (2)
		(n2) edge [loop above] node {$c$} (n2);		
	\end{tikzpicture}
	\caption[Sequence of most complex languages]{Automata from~\cite{DBLP:journals/ijfcs/Brzozowski13} yielding a sequence of most complex languages in the class of all languages. The sole final state is marked by a double line.}
	\label{syncmax:fig:most-complex-all-lang}
\end{figure}

As the full transformation semigroup is completely reachable, completely distinguishable,
and contains the full symmetric group, which is primitive (see \crefitem{sync_max:ex:primitivegrps}{sync_max:ex:primitivegrps:2}),
we obtain the following for the class of all regular languages.

\begin{proposition}
	Let $\mathcal A_n$ be a sequence of $n$-state minimal automata 
	such that $L_n = L(\mathcal A_n)$ is most complex in the class of all regular languages.
	Then $\TS{\mathcal A_n}$ equals the full transformation semigroup on $n$ states and so $\mathcal A_n$
	admits no non-trivial $\TS{\mathcal A_n}$-congruence
    and is completely reachable, completely distinguishable,
	and sync-maximal.	
\end{proposition}

So, for the class of all regular languages, we also obtain 
sync-maximal automata not admitting any non-trivial $\TS{\mathcal A}$-congruences.
In general, how do both properties, sync-maximality (or more generally,
the size of accepting automata for sets of synchronizing words, i.e., the operation $\mathcal A \mapsto \Syn(\mathcal A)$ that associates a regular ideal language with every automaton or regular language by using the minimal automaton) 
and absence of non-trivial $\TS{\mathcal A}$-congruences of the minimal automata, relate
to the other extremal properties of most complex languages? 

\begin{figure}[ht]
	\centering 	
\begin{tikzpicture}
	[node distance=5.5em and 5.8em, on grid, auto, state/.style = {draw = black, rounded rectangle, minimum width=4em, minimum height=1em},every loop/.style={looseness=10}]
	
	\node[state, initial, initial where=below] (1) {$1$};
	\node[state] (2) [right=of 1] {$2$};
	\node[state] (3) [xshift=0.5cm,right=of 2] {$3$};
	\node[state] (4) [right=of 3] {$4$};
	\node (dots) [right=of 4] {$\ldots$};
	\node[state] (n2) [right=of dots] {$n-2$};
	\node[state] (n1) [right=of n2] {$n-1$};
	\node[state, accepting] (n) [above=of 2] {$n$};
	
	\path[>=stealth]
	(1) edge[loop above] node[above] {$a,b,c,d,f$} ()
	(2) edge[loop below] node[below] {$c,d,e$} ()
	(3) edge[loop above] node[above] {$c,d,f$} ()
	(4) edge[loop above] node[above] {$b,c,d,f$} ()
	(n2) edge[loop above] node[above] {$b,c,d,f$} ()
	(n1) edge[loop above] node[above] {$b,f$} ()
	(n) edge[loop above] node[above] {$a,b,c,d,e,f$} ();
	
	\path[>=stealth]
	(1) edge node {$e$} (2)
	(2) edge node {$f$} (n)
	(2) edge[bend left=25] node[above] {$a,b$} (3)
	(3) edge[bend left=25] node[above] {$b,e$} (2)
	(3) edge node {$a$} (4)
	(4) edge[bend left=25] node[below] {$e$} (2)
	(4) edge node {$a$} (dots)
	(dots) edge node {$a$} (n2)
	(n2) edge node {$a$} (n1)
	(n2) edge[bend left=35] node[above] {$e$} (2)
	(n1) edge[bend right=50] node[above] {$a,c,e$} (2)
	(n1) edge[bend left=42] node[above] {$d$} (1);
\end{tikzpicture}
	\caption[Sequence of most complex languages w.r.t.~regular ideal languages]{Automata $\mathcal D_n$ from~\cite{DBLP:journals/dmtcs/BrzozowskiDL16} yielding a sequence of most complex languages in the class of regular ideal languages. The sole final state is marked by a double line.}
	\label{syncmax:fig:most-complex-ideals}
\end{figure}

For regular ideal languages, by \cref{prelim:eqn:reg_ideal},
there is no blow-up of minimal automata for 
the set of synchronizing words. Thus,
for a most complex sequence in the class of ideal languages
accepting automata are not sync-maximal.
For this class, Brzozowski~\cite{DBLP:journals/dmtcs/BrzozowskiDL16} 
gave, for $n \ge 5$, 
the sequence 
\[
\automat{D_n}{\{1,2,\ldots,n\}}{\{a,b,c,d,e,f\}}{\delta_n}{1}{\{n\}}
\] 
of minimal automata accepting most complex regular ideal languages where 
\begin{enumerate}
	\item $a$ induces the cyclic permutation $(2 ~ 3 ~ \cdots ~ n - 1)$,
	
	\item $b$ induces the transposition $(2 ~3)$, 
	
	\item  $c$ is a simple idempotent letter mapping $n - 1$ to $2$, 
	
	\item  $d$ is a simple idempotent letter mapping $n - 1$ to $1$,
	
	\item $e$ maps everything in $\{1,\ldots,n-1\}$ to $2$ and fixes $n$,
	
	\item $f$ is a simple idempotent letter mapping $2$ to $n$.
\end{enumerate}
See also \cref{syncmax:fig:most-complex-ideals} for a state diagram.
Here, $|\TS{\mathcal D_n}| = n^{n-2} + (n-2)2^{n-2} + 1$ (the proof
of this result appears in~\cite{BrzozowskiSY18})
and a minimal automaton for reversal of $L((\mathcal D_n)_{\{a,d,e,f\}})$ 
has size $2^{n-2} + 1$~\cite{DBLP:journals/dmtcs/BrzozowskiDL16}.
So, w.r.t.~other properties than the state complexity of $\Syn(\mathcal D_n)$
we obtain an exponential blow-up. Next, in \cref{sync_max:prop:Dn},
we show that these automata do not admit any non-trivial $\TS{\mathcal D_n}$-congruence
and that an automaton constructed from the subsets
with a property similar to distinguishability
indeed is minimal and exponentially larger.

\begin{proposition}\label{sync_max:prop:Dn}
\begin{enumerate}
\item \label{sync_max:prop:Dn:1} The automata $\mathcal D_n$ admit no non-trivial $\TS{\mathcal D_n}$-congruence.

\item \label{sync_max:prop:Dn:item:reach} Every subset $A \subseteq \{1,2,\ldots,n\}$ with $\{1,n\} \subseteq A$ is reachable.

\item \label{sync_max:prop:Dn:item:dist} For distinct $A,B \subseteq \{1,2,\ldots,n\}$ with $\{ 1,n \} \subseteq A \cap B$
 there exists $u \in \{a,b,c\}^*f$ mapping precisely one subset to  $\{1,n\}$
 and the other to a set containing at least three states.
 
\item \label{sync_max:prop:Dn:4} The minimal automaton for $L(\mathcal P(\mathcal D_n)_{\{1,2,\ldots,n\}, \{ \{1,n\}\}})$
 has at least $2^{n-2}$ states.
\end{enumerate}
\end{proposition}
\begin{proof} 	Set $S = \{2,3,\ldots,n-1\}$.
	Since $(2 ~ 3 ~ \cdots ~ n - 1)$ and $(2~3)$ restricted to the domain~$S$
	generate the full symmetric group on~$S$,
	$\TM{\restrictAutGamS{(\mathcal D_n)}{\{a,b\}}{S}}$
	contains every permutation of~$S$ that fixes~$1$.
	As $c$ is a simple idempotent when restricted to $S$
	and the symmetric group is primitive,
	by~\cref{sync_max:thm:prim},
	$\restrictAutGamS{(\mathcal D_n)}{\{a,b,c\}}{S}$
	is completely reachable and completely distinguishable (note that $\restrictAutGamS{(\mathcal D_n)}{\{a,b,c\}}{S}$
	is isomorphic to the automaton from \cref{syncmax:fig:most-complex-all-lang}.)
	
\ref{sync_max:prop:Dn:1}:  Let  $\sim$ be a $\TM{\mathcal D_n}$-congruence 
with $p \sim q$ for states $p,q$ with $p < q$.    
There exists $u \in a^*f$
such that $\delta_n(q, u) = n$
and $\delta_n(p, u) \in \{1,2,\ldots,n-1\}$.
Now, for any $r \in \{1,\ldots,n-1\}$
we find $v \in \{a,d,e\}^*$
such that $\delta_n(p, uv) = r$.
Then $\delta_n(q, uv) = n$
and so $r \sim n$. Hence, $\sim$ equals the universal relation.

\ref{sync_max:prop:Dn:item:reach}: 	We have $\delta_n(\{1,2,\ldots,n\}, ea^{n-3}d) = \{1,n\}$.
Next, let $A \subseteq \{1,2,\ldots,n\}$ be a subset that
properly contains $\{1,n\}$.
As $\restrictAutGamS{(\mathcal D_n)}{\{a,b,c\}}{S}$ is completely reachable there exists $u \in \{a,b,c\}^*$
such that $\applytrans{\delta_n}{S}{u} = A \setminus \{1,n\}$.
Since $u$ fixes $1$ and $n$, we find $\applytrans{\delta_n}{\{1,2,\ldots,n\}}{u} = A \cup \{1,n\} = A$.

\ref{sync_max:prop:Dn:item:dist}: Consider $A \cap S$ and $B \cap S$, which can be assumed to be non-empty, 
since otherwise exactly one of them equals $\{1,n\}$.
As $\restrictAutGamS{(\mathcal D_n)}{\{a,b,c\}}{S}$ is completely distinguishable,
there exists $u \in \{a,b,c\}^*$ mapping precisely
one subset, say $A \cap S$, to a singleton.
As $a$ cyclically permutes $S$, we can assume $\delta_n(A \cap S) = \{2\}$.
Then $\delta_n(A, uf) = \delta_n(\{1,n,2\},f) = \{1,n\}$ 
and $|\delta_n(B, uf)| \ge 3$
since 
$\delta_n(B, uf) = \delta_n(\{1,n\}, uf) \cup \delta_n(B \cap S, uf)$
where $|\delta_n(B \cap S, u)| \ge 2$
and $\delta_n(\{1,n\}, u) = \{1,n\}$.

\ref{sync_max:prop:Dn:4}: There are $2^{n-2}$ sets of the form $\{1,n\} \cup A$
with $A \subseteq S$,
each reachable by~\ref{sync_max:prop:Dn:item:reach} and, by~\ref{sync_max:prop:Dn:item:dist}, every
pair of distinct sets among these sets is 
distinguishable in $\mathcal P(\mathcal D_n)_{\{1,2,\ldots,n\}, \{\{1,n\}\}}$.
So, \crefitem{prelim:thm:char_minaut}{thm:minaut:item:statesdist}
gives the claim.
\end{proof}

Note that $ef \in \Syn(\mathcal D_n)$. While, by \cref{prelim:eqn:reg_ideal},
the automata $\mathcal D_n$ themselves accept $\Syn(\mathcal D_n)$,
\Cref{sync_max:prop:Dn} shows that we obtain exponentially larger
minimal automata w.r.t.~different final state sets in the power automata
whose start states are the whole state set.
So, in view of the observations made in \cref{sync_max:prop:Dn},
we ask the following questions.

\begin{center}
	\fbox{ \begin{minipage}{0.8\linewidth}
			How do families of most complex regular languages $L_n = L(\mathcal A_n)$ in some class and with minimal 
			automata $\mathcal A_n$  
			relate to $\TS{\mathcal A_n}$-congruences?
		    Furthermore, in a family of most complex regular languages w.r.t.~some class,
		    can we choose final state sets in the power automata with start state the whole state set to obtain
		    exponentially larger minimal automata; and conversely
		    how does existence of those final state sets relate
		    to extremal properties w.r.t.~the complexity measures
		    of most complex families of regular languages?
		\end{minipage}
	}
\end{center}

We mention in passing that it has been observed
that state-wise maximal automata
for the results of a wide class of regularity-preserving
operations can often be realized by so-called \emph{monsters}~\cite{DBLP:journals/dmtcs/CaronCLP20,DBLP:phd/hal/court20}. 
These are automata that contain a letter for each transformation
of the underlying state set (and so preserve no congruence, are completely reachable, and are completely distinguishable). However, there is no notion of a monster relative to some class of languages. 

\paragraph{Primitivity and $A$-Primitivity.}
Most characterizations of primitivity
from \cref{sec:prim_groups}
transfer easily to the setting of strongly $\{a,b\}$-primitive
groups. In this sense, the properties behind these characterizations
are actually ``local'' properties applying to pairs
of distinct points, yielding primitivity when using them ``globally'', i.e.,
applying them to all points or small subsets of points,
via \cref{sync_max:prop:primitive_Aprimitive}.
However, the characterizations
\crefitem{sync_max:thm:prim_char:prim}{sync_max:thm:prim_char2:strongly_factorizable}
and \cref{sync_max:thm:prim_known_char}
do not seem to admit a (at least straightforward) generalization
to strongly $\{a,b\}$-primitive groups.
\Cref{sync_aut:permgrp:ex:Aprim}
gave a strongly $\{0,1\}$-primitive group $G$
that is imprimitive, and so where no 
point stabilizer is maximal.
Furthermore, for $t \colon [8] \to [8]$
with $t(0) = t(2) = t(4) = t(6) = 0$
and $t(1) = t(3) = t(5) = t(7) = 1$
there is no $g \in G$
such that $tg$ is idempotent.
This implies that $\langle G \cup \{t\} \rangle_S$
cannot be written in the form $EG$
for its set of idempotents~$E$.
In this sense, the characterizations by maximal point stabilizers
and the one on adjoining transformations of rank two
are genuinely ''global`` properties characterizing primitivity
involving all points, or their point stabilizers,
and it breaks down when only
considering a subset $A$ of points (or the point stabilizers of points from $A$)
in imprimitive permutation groups that are $A$-primitive. 




\paragraph{Permutation Groups and Non-Permutations.} Pin~\cite{DBLP:conf/icalp/Pin78} 
proved that a transitive permutation group of prime degree
synchronizes any non-permutation (which follows
from~\cite[Lem.~3.1 \& Cor.~3.3]{ArnoldS06} as well).
This result was also stated by Salomaa~\cite[Thm.~5]{DBLP:journals/jucs/Salomaa02} (without reference to Pin).
In~\cite{DBLP:journals/eatcs/Salomaa00,DBLP:journals/jucs/Salomaa02,DBLP:journals/tcs/Salomaa03,DBLP:conf/birthday/Salomaa12},
besides Pin's result, 
he gave further sufficient conditions on sets of functions to generate
a given function, or a function of a given rank (where synchronizability corresponds to functions of rank one).
For example, he also stated 
that a $2$-transitive group synchronizes every non-permutation~\cite[Thm.~4]{DBLP:journals/jucs/Salomaa02} (this follows from \crefitem{sync_max:prop:sync-between}{sync_max:prop:sync-between:2}).
He also investigated  the length of generating sequences, i.e., words over the generators, that realize functions
from some class.
The focus of these studies was not on algorithmic questions, but 
the problems investigated also overlap with the algorithmic problems
mentioned in \cref{prelim:compl:sec:discussion}, in particular
those related to transformation semigroups discussed at the end.
In a broader context, Salomaa was motivated by questions in logic, and under what conditions a set of basic logical functions
generate all logical functions.
In~\cite{Salomaa60,Salomaa1964OnBG} 
he investigated sets of functions  $f \colon [n]^i \to [n]$, $i \ge 0$, generating all $m$-ary functions
for every $m \ge 0$ under substitutions when adding a surjective function $[n]^j \to [n]$, $j \ge 2$,
depending on at least two input variables.
At the intersection between group theory and semigroup theory, one of the driving forces is the paradigm
of ``detecting properties [of permutation groups] with functions''~\cite{ArCamSt2017}, or, 
as formulated in~\cite{ARAUJO2021a,araujo2016}, the following related question:
\begin{quote}
	Classify the pairs $(G, a)$, where $a$ is a map on a finite set $\Omega$ and $G$
	is a group of permutations on $\Omega$, such that the semigroup
	$\langle G, a \rangle$ generated by $a$ and~$G$ has a given property $P$.
\end{quote}
The results of the present chapter fit into this paradigm.
Other results are given in~\cite[Sec.~8]{ArCamSt2017}, as
for example, that $2$-homogeneity is characterized by the property~$P$
of containing every non-permutation~\cite[Thm.~3.18]{mcalister1998semigroups} (see also~\cite[Thm.~8.2]{ArCamSt2017}).

For further connections, or similar properties, we refer to~\cite{andre2016,araujo2019,araujo2021,araujo2013,AraujoCameron2014,ArCamSt2017,doi:10.1112/plms/pdw040,bray2020}.
Göhring~\cite{DBLP:journals/jalc/Gohring97}
has investigated so-called \emph{\v{C}ern\'y-like automata}
by various properties of permutation groups contained in the transformation
semigroup of an automaton.

For the result from \cref{sync_max:sec:perm_grps}
it does not matter whether we formulate them in terms of the generated 
semigroup or monoid, as for non-empty $P \subseteq \mathcal S_n$
and $T \subseteq \mathcal T_n \setminus \mathcal S_n$ we have 
$
\langle P \cup T \rangle_S = \langle P \cup T \rangle_M
$
as $\operatorname{id}_{[n]} \in \langle P \rangle_S$.




% Lässt sich nicht kompilieren ohne Fehler. Auch Fehler, wenn ich es in den "note"-Teil der bib-Datei einbinde
% \footnote{The original two parts of Lagranges work
	%	are also available as google scans (accessed at 2026-01-04): 
	%	\url{https://books.google.de/books?id=NQQOAAAAQAAJ&hl=de&pg=RA1-PA134#v=onepage&q&f=false}
	%	and \url{https://books.google.de/books?id=enPaeQG_cfwC&hl=de&pg=RA1-PA138#v=onepage&q&f=false}.}	
\IN{O'Nan-Scott theorem}
\IN{socle of a group}
\IN{twisted wreath product action}
\paragraph{Historical Remarks on Primitive Permutation Groups.} %\'Evariste 
Galois~\cite{Neumann2011} introduced primitive permutation groups  
in his seminal approach
 to the characterization
of the solvability of polynomial equations,
based on an observation made by Lagrange~\cite{lagrangecollected2}
relating permutation groups to polynomial equations.\footnote{See also \url{https://hsm.stackexchange.com/questions/15489/where-did-lagrange-write-his-technique-of-resolvents-for-solving-polynomials?noredirect=1&lq=1} for links to the original sources.}
The first systematic treatment of permutation group theory 
was written by 
Jordan~\cite{jordan1870},
where he proved many, now classical, results about primitive permutation groups
as for example \cref{sync_max:thm:prim_char:jordan}.
An intransitive permutation group can be decomposed as a subcartesian product of its transitive constituents~\cite[Thm.~1.2]{cameron_1999}.
Now, a transitive but imprimitive permutation group
can be decomposed into 
a wreath product 
where the action of the group equals the imprimitive
action of the wreath product for an identification
of the permutation domain with a Cartesian product,  
see~\cite[Thm.~1.8]{cameron_1999},~\cite{kovacs1989}, \cite[Sect.~2,6]{dixon1996permutation},
or~\cite[Sect.~4.4]{Kurzweil04} for a very detailed explanation.
 
Scott~\cite[p.~329]{Scott80} 
classified the subgroups of the symmetric group up to conjugation
(with a footnote mentioning that O'Nan obtained
the same result independently). The cases fall into categories covering the primitive and imprimitive groups separately (which can be seen from the proof).
So, the result yields a classification
of the primitive permutation groups, called the \emph{O'Nan-Scott theorem} 
and explicitly given by Cameron~\cite[Thm.~4.1]{cameron81} via the action
on the \emph{socle}, i.e., the subgroup generated by the minimal normal subgroups (see~\cite[Sect.~8]{cameron81} for a discussion in a broader context). 
The cases from the statement 
involving $2$-transitive groups, which were either almost simple or affine,
were given previously by Burnside~\cite{burnside1911}. However, the result was even given
earlier by Jordan~\cite{jordan1870} as pointed out by Peter Neumann, 
see the blog post~\cite{cameron:blog}.
However, a case was missing in the statement of the theorem, which was discovered 
and corrected by Aschbacher \& Scott~\cite{ASCHBACHER1985} and, independently, Kov\'acs~\cite{KOVACS1986114}.
Liebeck et al.~\cite{Liebeck88} gave a fully self-contained
and corrected statement of the O'Nan-Scott theorem, where the previously missing
case was described using the \emph{twisted wreath product action}.
For further information, we refer to the surveys~\cite{cameron81,praeger90}
and the textbooks~\cite{cameron_1999,dixon1996permutation}.